

\section{The $(d,d^{\prime})$-Dimensional Toric Code}
The main error correction code that we will use is the Toric code invented by Kitaev. In the 2D case, the code is obtained by taking a Toric, tiling it with squares, setting qubits on the edges, and checks as follows: For every vertex, the parity of the qubits adjoining it is even, and for any square, the parity of the qubits lying on its sides, in the phase basis, is even. A high-dimensional version can be obtained by taking a high-dimensional Toric, encoded as a cubical complex, and then setting the $X$-checks, the qubits, and the $Z$-checks to be faces at sequenced dimensions. Throughout this paper, qubits will be set on $d^{\prime}$-dimensional faces, and the $X$-checks and $Z$-checks on the $d^{\prime} \mp 1$ dimensional faces. For a review, we refer to \cite{Kitaev1997QuantumCA}, \cite{Dennis_2002}.

For defining a decoder and proving its 'correctness,' we will have to overload notation. We derive the Toric complex from the skeleton of the Cayley graph of the cyclic group, using the generating subset as a collection of positive directions. The $\rotfvs{}$ encodes a face, one vertex $v$ that serves as its 'root' and a subset of directions $S^{\prime}$. Formally: 

\begin{definition}[$d$-Dimensional Toric Complex] 
  \label{def:toric_complex}
  For integers $d$ and $n$, denote the group resulting from the $d$-directed power of the cyclic group of order $n$ by $\mathcal{G}_{n}^{d} = \mathbb{Z}_{n}^{\times d}$, and denote the standard generating set of $\mathcal{G}_{n}^{d}$ by $S = \{ 1_{i} \colon i \in [d] \}$. We define the \textit{$d$-dimensional Toric complex} to be the hypergraph obtained from the Cayley graph $\text{Cay}(\mathcal{G}_{n}^{d}, S)$ by taking its vertices as the $0$-faces and for any $i \in [d]$ we define the $i$-faces as follow: First, for a vertex $v \in \mathcal{G}_{n}^{d}$ and for a subset of generators $S^{\prime} \subset S$ of size $|S^{\prime}| = i$, the $i$-face rooted at $v$ and spanned by $S^{\prime}$, denoted by $\theta_{v, S^{\prime}}$, is:
  \begin{equation*}
    \begin{split}
      \theta_{v,S^{\prime}} \colonequals \bigcup_{I \subset S^{\prime}} \left\{  \left(\prod_{g \in I} g \right) v  \right\}
    \end{split}
  \end{equation*}
  The $i$-faces of the complex are the collection of all the possible rooted spanned $i$-faces induced by $\text{Cay}(\mathcal{G}_{n}^{d})$. We will denote them by $\mathcal{F}_{i}$. 

  Additionally, we name $n$ and $d$, the \textit{side length} and the \textit{dimension} of the complex.  
\end{definition}

Next, we formalize binary assignments on the Toric complex and introduce boundary operators. These boundary operators are the same as standard boundary operators in the chain complex formalism.

\begin{definition}[Assignments on a Toric complex]
  \label{def:asontoric}
  Let $\mathcal{T}$ be a $d$-dimensional Toric complex. For each $i$, the set of $i$-faces induces a vector space $\mathbb{F}_{2}^{|\mathcal{F}_{i}|}$. An element $z\in \mathbb{F}_{2}^{|\mathcal{F}_{i}|}$ is called an \textit{assignment on the $i$-faces} of $\mathcal{T}$.

  For a vertex $v$ and a subset of generators $S^{\prime}\subset S$ of size $|S^{\prime}|=i$, let $\rotfvs{}$ denote the standard basis vector corresponding to the $i$-face $\rotfs{}$. Thus the family
  \begin{equation*}
    \begin{split}
      \bigl\{ \rotfv{v}{S^{\prime}} : v\in\mathcal{F}_{0},\ S^{\prime}\subset S,\ |S^{\prime}|=i \bigr\}
    \end{split}
  \end{equation*}
  is the standard coordinate basis of $\mathbb{F}_{2}^{|\mathcal{F}_{i}|}$.

  Let $z$ be an assignment on the $i$-faces. Let $v_{1}, \ldots, v_{k} \in \mathcal{F}_{0}$ and $S^{\prime}_{1}, \ldots, S^{\prime}_{k} \subset S$ be such that $\theta_{v_{1},S^{\prime}_{1}}, \theta_{v_{2},S^{\prime}_{2}}, \ldots, \theta_{v_{k},S^{\prime}_{k}}$ are precisely the $i$-faces in the support of $z$, namely:
  \begin{equation*}
    \begin{split}
    z = \sum_{i=1}^{k}\rotfv{v_{i}}{S^{\prime}_{i}}
    \end{split}
  \end{equation*} 
  We refer to $\theta_{v_{1},S^{\prime}_{1}}, \theta_{v_{2},S^{\prime}_{2}}, \ldots, \theta_{v_{k},S^{\prime}_{k}}$ as the collection of faces that form $z$, or briefly, the faces form $z$.
  For $i > 0$, we define the linear map $\partial^{-}: \mathbb{F}_{2}^{|\mathcal{F}_{i}|} \rightarrow \mathbb{F}_{2}^{|\mathcal{F}_{i-1}|}$ such that
  \begin{equation*}
    \begin{split}
      \partial^{-} \rotfvs{} \colonequals \sum_{g \in S^{\prime}} \left( \rotfv{v}{S^{\prime}/g} + \rotfv{gv}{S^{\prime}/g} \right)
    \end{split}
  \end{equation*}
  For $i < d$, we define the linear map $\partial^{+}: \mathbb{F}_{2}^{|\mathcal{F}_{i}|} \rightarrow \mathbb{F}_{2}^{|\mathcal{F}_{i+1}|}$ such that
  \begin{equation*}
    \begin{split}
      \partial^{+} \rotfvs{} \colonequals \sum_{g \in S/S^{\prime}} \left( \rotfv{v}{S^{\prime} \cup g} + \rotfv{g^{-1}v}{S^{\prime}\cup g} \right) 
    \end{split}
  \end{equation*}
  We define the \textit{restriction of $z$ at vertex $v$} to be the linear map $|_{v} \colon \mathbb{F}_{2}^{|\mathcal{F}_{i}|} \rightarrow \mathbb{F}_{2}^{|\mathcal{F}_{i}|}$ such that: 
  \begin{equation*}
    \begin{split}
      \rotfv{u}{S^{\prime}} |_{v} \colonequals    
      \begin{cases}
      \rotfv{u}{S^{\prime}}  & v \in \theta_{u,S^{\prime}}  \\ 
        0 & \text{ else }
      \end{cases}
    \end{split}
  \end{equation*}
  We refer to the faces rooted at $v$ as the positive faces of $v$, and define the \textit{positive restriction at vertex $v$} to be the linear map $|_{v+} \colon \mathbb{F}_{2}^{|\mathcal{F}_{i}|} \rightarrow \mathbb{F}_{2}^{|\mathcal{F}_{i}|}$ such that: 
  \begin{equation*}
    \begin{split}
      \rotfv{u}{S^{\prime}} |_{v+} \colonequals    
      \begin{cases}
      \rotfv{u}{S^{\prime}}  & u = v \\ 
        0 & \text{ else }
      \end{cases}
    \end{split}
  \end{equation*}
  

  Two $i$th dimensional faces will be said to be connected if their intersection is $i-1$th dimensional face. So, an assignment $z$ is \textit{connected} if the graph that is induced by taking the faces that support $z$ as vertices and defining the edges to be the adjacency relation between them is connected. We define the connected components of $z$ to be its subsets such that each of them is a connected assignment. For a $j$-face $f$, we say that $f \in z$ if $z$ has a support on $\rotfvs{}$ such that $f \in \rotfs{}$. Consider an assignment $z \in \mathbb{F}_{2}^{\mathcal{F}_{i}}$ and a vertex $v \in \mathcal{G}_{n}^{d}$. We say that $z$ is a \textit{rooted assignment} at vertex $v$ if there are subsets $S^{\prime}_{1}, \ldots, S^{\prime}_{k} \subset S^{\times k}$ such that each of $S^{\prime}_{j}$ is at size $i$ and $z$ is decomposed as:
  \begin{equation*}
    \begin{split}
    z = \sum_{j}^{k}\rotfv{v}{S^{\prime}_{j}}
    \end{split}
  \end{equation*}
  Put differently, $z$ is the result of taking a collection of $i$-faces that are rooted at $v$.
\end{definition}


%We will use only this CSS property and the local geometry of the code; standard facts about distance and rate may be found in the toric-code literature, for example in \cite{Dennis_2002}.
We are ready to give a formal definition of the $(d,d^{\prime})$-Toric code. 


\begin{definition}[$(d,d^{\prime})$ Toric Code] 
  Let $d, d^{\prime}$ be integers such that $1 \le d^{\prime} < d$. For any integer $n$, we construct the $n$th instance of the quantum code family $(d,d^{\prime})$ \textit{Toric Code} as follows: Let $G = (\mathcal{F}_{0}, \ldots, \mathcal{F}_{d})$ be the $d$-dimensional Toric complex obtained from $\mathcal{G}_{n}^{d}$ as defined in \Cref{def:toric_complex}. We associate the (q)bits with the $\mathcal{F}_{d^{\prime}}$, the $X$-checks with the $\mathcal{F}_{d^{\prime}-1}$, and the $Z$-checks with the $\mathcal{F}_{d^{\prime}+1}$. An $X$-check, $c_{x}$, is satisfied if the parity of (q)bits set on the $d^{\prime}$-face containing $c_{x}$ is even. Similarly, a $Z$-check, $c_{z}$, is satisfied if the parity, in the Hadamard basis, of the (q)bits set on the $d^{\prime}$-faces contained in $c_{z}$ is even.
\end{definition}

We argue that the $(d,d^{\prime})$-Toric code is indeed a CSS, and prove that by showing that boundary operators $\partial^{\pm} :\mathbb{F}_{2}^{|\mathcal{F}_{i}|} \rightarrow \mathbb{F}_{2}^{|\mathcal{F}_{i \pm 1}|}$ induce a chain complex.
\begin{claim}[CSS structure]
  \label{claim:csscode}
  The $(d,d^{\prime})$-Toric code is a CSS code.
\end{claim}
\begin{proof}
  We will show that $ \mathbb{F}_{2}^{|\mathcal{F}_{i}|} \xrightarrow{ \partial^{\pm} } \mathbb{F}_{2}^{|\mathcal{F}_{i \pm 1}|}$ are a chain complex and its dual chain. First, we show that $\left( \partial^{+} \right)^\top = \partial^{-}$. Consider vertices $v, u \in V$ and two generator subsets $A,B \subset S$. Consider the basis elements defined by them: $\rotfv{v}{A}$ and $\rotfv{u}{B}$. Suppose that:
  \begin{equation*}
    \begin{split}
       ( \partial^{+} \rotfv{v}{A} ) \rotfv{u}{B}   = 1 
    \end{split}
  \end{equation*}
Means that there exists $g \notin A$ such that $A \cup \{g\} = B$ and either $v = u$ or $g^{-1}v = u$. But that is equivalent to saying that there is $g \in B$ such that $B \setminus \{g\} = A$ and either $v = u$ or $v = g^{-1}u$. In other words:
  \begin{equation*}
    \begin{split}
       & ( \partial^{+} \rotfv{v}{A} ) \rotfv{u}{B}   =   \rotfv{v}{A} ( \partial^{-} \rotfv{u}{B} ) \\ 
       \Rightarrow & \left( \partial^{+} \right)^\top = \partial^{-} 
    \end{split}
  \end{equation*}

  So it's left to show that the boundary of a boundary is zero. Consider a basis vector $\rotfvs{}$:
  \begin{equation*}
    \begin{split}
      \partial^{-} \left(   \partial^{-} \rotfvs{} \right) & = \partial^{-} \left( \sum_{g \in S^{\prime}} \rotfv{v}{S^{\prime}/g} + \sum_{g \in S^{\prime}} \rotfv{gv}{S^{\prime}/g}  \right) \\
      & = \sum_{ \substack{ g^{\prime},g \in S^{\prime} \\ g^{\prime}\neq g } } \left( \rotfv{v}{S^{\prime}/\{ g,g^{\prime} \} } +   \rotfv{g^{\prime}v}{S^{\prime}/\{ g,g^{\prime} \} }+ \rotfv{gv}{S^{\prime}/\{ g,g^{\prime} \} } + \rotfv{g^{\prime}gv}{S^{\prime}/\{ g, g^{\prime}\}} \right) \\
      & = 0
    \end{split}
  \end{equation*}
Where we used that the term in the closure is invariant to sweeping between $g$ and $g^\prime$, and therefore any term in the sum appears twice. Similarly:
  \begin{equation*}
    \begin{split}
      \partial^{+} \left(   \partial^{+} \rotfvs{} \right) & = \partial^{+} \left( \sum_{g \in S/S^{\prime}} \rotfv{v}{S^{\prime} \cup g} + \sum_{g \in S / S^{\prime}} \rotfv{g^{-1}v}{S^{\prime} \cup g}  \right) \\
      & = \sum_{ \substack{ g^{\prime},g \in S/S^{\prime} \\ g^{\prime}\neq g } } \left( \rotfv{v}{S^{\prime}\cup\{ g,g^{\prime} \} } +   \rotfv{g^{\prime -1}v}{S^{\prime}\cup\{ g,g^{\prime} \} }+ \rotfv{g^{-1}v}{S^{\prime}\cup\{ g,g^{\prime} \} } + \rotfv{g^{\prime-1}g^{-1}v}{S^{\prime}\cup\{ g, g^{\prime}\}} \right) \\
      & = 0
    \end{split}
  \end{equation*}

The conditions $ (\partial^{\pm})^{\circ 2} = 0 $ imply that $ \mathbb{F}_{2}^{|\mathcal{F}_{i}|} \xrightarrow{ \partial^{\pm} } \mathbb{F}_{2}^{|\mathcal{F}_{i \pm 1}|}$ are chain complexes. Therefore, for any $d^{\prime}$, every $d^{\prime}-1$ face and every $d^{\prime}+1$ face overlap on an even number of $d^{\prime}$-faces, so the corresponding $X$- and $Z$-checks commute. This is exactly the CSS condition.
\end{proof}

 
\begin{definition}[Positive Neighborhood]
  \label{def:positiveNei}
For a vertex $v$ and a chain $z$, we use the notation $\textbf{Str}_{0+}(v,z)$ to denote the vertices in $z$ that are connected to $v$ through a stepping via generator, namely all the vertices that support a face in $z$ of the form $gv$ for some $g \in S$. We think of $\textbf{Str}_{0+}(v,z)$ as the positive neighborhood of $v$ in $z$.
\end{definition}

